% Equation sheet of the nacre-shell screen (session P0-W3.1).
% Symbolic only: no literature value appears here. Values are fetched inputs that enter
% the implementation through the parameter table; paper text quotes numbers only through
% generated macros. The symbol table is generated from theory/symbols.yaml.
% Build: latexmk -pdf -outdir=theory/build theory/equation_sheet.tex (from the repository root,
% or `latexmk -pdf -outdir=build equation_sheet.tex` inside theory/).
\documentclass[9pt]{extarticle}
\usepackage[a4paper,margin=11mm,top=10mm,bottom=12mm]{geometry}
\usepackage{amsmath}
\usepackage{array}
\usepackage{multicol}
% Notation shared by theory/equation_sheet.tex and theory/derivation_body.tex.
% Every macro here corresponds to an id in theory/symbols.yaml. Requires amsmath.
\providecommand{\kB}{k_\mathrm{B}}
\providecommand{\NA}{N_\mathrm{A}}
\providecommand{\Rg}{R}
\providecommand{\Fc}{F}
\providecommand{\dd}{\mathrm{d}}
\providecommand{\e}{\mathrm{e}}
% interfacial free energies
\providecommand{\gow}{\gamma_\mathrm{ow}}
\providecommand{\gpw}{\gamma_\mathrm{pw}}
\providecommand{\gop}{\gamma_\mathrm{op}}
\providecommand{\go}{\gamma_\mathrm{o}}
\providecommand{\gp}{\gamma_\mathrm{p}}
\providecommand{\gw}{\gamma_\mathrm{w}}
\providecommand{\gnl}{\gamma_\mathrm{nl}}
\providecommand{\gsl}{\gamma_\mathrm{sl}}
\providecommand{\gslz}{\gamma_\mathrm{sl}^{(0)}}
\providecommand{\dga}{\Delta\gamma_\mathrm{ads}}
\providecommand{\gsn}{\gamma_\mathrm{sn}}
\providecommand{\gsnz}{\gamma_\mathrm{sn}^{(0)}}
\providecommand{\dgw}{\Delta\gamma_\mathrm{w}}
\providecommand{\geff}{\gamma_\mathrm{eff}}
\providecommand{\Wwet}{W_\mathrm{adh,wet}}
\providecommand{\Wdry}{W_\mathrm{adh,dry}}
% sequestration and supersaturation
\providecommand{\DGseq}{\Delta G_\mathrm{seq}}
\providecommand{\DGsite}{\Delta G_\mathrm{site}}
\providecommand{\Eenr}{\mathcal{E}}
\providecommand{\cb}{c_\mathrm{bulk}}
\providecommand{\cl}{c_\mathrm{local}}
\providecommand{\cmax}{c_\mathrm{max}}
\providecommand{\Gs}{\Gamma_\mathrm{s}}
\providecommand{\ds}{\delta_\mathrm{s}}
\providecommand{\Vsite}{V_\mathrm{site}}
\providecommand{\GCa}{\Gamma_\mathrm{Ca}}
\providecommand{\GCasl}{\Gamma_\mathrm{Ca,sl}}
\providecommand{\GCasn}{\Gamma_\mathrm{Ca,sn}}
\providecommand{\aCa}{a_{\mathrm{Ca}^{2+}}}
\providecommand{\aCO}{a_{\mathrm{CO}_3^{2-}}}
\providecommand{\Ksp}{K_\mathrm{sp}}
\providecommand{\Om}{\Omega}
\providecommand{\Omb}{\Omega_\mathrm{bulk}}
\providecommand{\Oml}{\Omega_\mathrm{local}}
\providecommand{\Dmu}{\Delta\mu}
% nucleation and morphology
\providecommand{\Vm}{V_\mathrm{m}}
\providecommand{\vm}{v_\mathrm{m}}
\providecommand{\DGhom}{\Delta G^{*}_\mathrm{hom}}
\providecommand{\DGhet}{\Delta G^{*}_\mathrm{het}}
\providecommand{\Ssp}{S_\mathrm{spread}}
\providecommand{\Ns}{N_\mathrm{s}}
\providecommand{\hc}{h_\mathrm{c}}
\providecommand{\hcoal}{h_\mathrm{coal}}
% stability
\providecommand{\ESHE}{E_\mathrm{SHE}}
\providecommand{\ELi}{E_\mathrm{Li}}
\providecommand{\muLi}{\mu_\mathrm{Li}}
\providecommand{\muLiz}{\mu_\mathrm{Li}^{\circ}}
\providecommand{\DGpbx}{\Delta G_\mathrm{pbx}}
% film
\providecommand{\rhoe}{\rho_\mathrm{e}}
\providecommand{\phiB}{\phi_\mathrm{B}}
\providecommand{\jreq}{j_\mathrm{req}}
% regimes and units
\providecommand{\Rs}{\textsf{R-slurry}}
\providecommand{\Rc}{\textsf{R-cell}}
\providecommand{\un}[1]{\quad[\mathrm{#1}]}

\setlength{\parindent}{0pt}
\setlength{\parskip}{2pt}
\setlength{\columnsep}{6mm}
\setlength{\abovedisplayskip}{3pt}\setlength{\belowdisplayskip}{3pt}
\setlength{\abovedisplayshortskip}{2pt}\setlength{\belowdisplayshortskip}{2pt}
\newcommand{\hd}[1]{\par\smallskip{\normalsize\bfseries #1}\par\nopagebreak}
\newcommand{\phys}[1]{\textit{#1}}
\pagestyle{empty}
\begin{document}

{\large\bfseries Equation sheet: governing relations of the CaCO$_3$-barrier screen}\par
{\small P0-W3.1. Symbolic only; every symbol is in the table overleaf and in \texttt{theory/symbols.yaml}; units in brackets.}

\smallskip
\begin{tabular}{@{}p{\dimexpr\textwidth*10/97\relax}p{\dimexpr\textwidth*41/97\relax}p{\dimexpr\textwidth*41/97\relax}@{}}
\hline
& \Rs: aqueous, alkaline, room temperature; casting and drying on native Al$_2$O$_3$/AlOOH
& \Rc: non-aqueous, fluorine-free electrolyte, collector at anodic potential vs Li$^+$/Li\\
\hline
Links that apply & C1 adhesion (1), C2 sequestration (2), supersaturation (3), C3 nucleation and Young (4, 5), template action (6), C4 morphology (7), C5 aqueous Pourbaix (8a)
& C5 Li-open stability window (8b), C4 insulating-film bound (9). Not defined here: (1)--(7), since there is no hydrated Ca$^{2+}$ reference and no aqueous $\Om$\\
\hline
\end{tabular}

\begin{multicols}{2}
\hd{1\quad Adhesion in water (C1, \Rs)}
\phys{The polymer never meets a dry oxide: every contact first displaces interfacial water, so the decisive work separates the contact into two \emph{water-wetted} surfaces.}
\begin{align}
\Wwet &= \gow + \gpw - \gop \un{J\,m^{-2}}\\
&= \Wdry + 2\gw - W_\mathrm{ow} - W_\mathrm{pw} \nonumber\\
\Wdry &= \go+\gp-\gop \un{J\,m^{-2}}\nonumber\\
\Wwet &= \Delta F_\mathrm{PMF}/A_\mathrm{c} \un{J\,m^{-2}}
\end{align}
Anchored against water iff $\Wwet>0$; $W_\mathrm{ow},W_\mathrm{pw}$ are the dry works of the oxide and polymer with water.

\hd{2\quad Sequestration (C2, \Rs)}
\phys{The template raises local Ca$^{2+}$ only if it holds Ca$^{2+}$ more strongly than water does; the cycle starts from hydrated Ca$^{2+}$ at the regime activity.}\par
Cycle: Ca$^{2+}$(aq, bulk, $\aCa$) $\to$ Ca$^{2+}$(binding layer). Solution side: molal standard state at the regime ionic strength; layer side: site-volume standard state, $\DGsite=-\Rg T\ln\!\big[\Vsite^{-1}\!\int_\mathrm{site}\e^{-w(\mathbf r)/\kB T}\dd^3r\big]$, with $w$ the per-ion PMF ($w\to0$ in bulk). With one coordinating group that must be deprotonated:
\begin{align}
\DGseq &= \DGsite + \Rg T\ln\!\big(1+10^{\,\mathrm{p}K_\mathrm{a}-\mathrm{pH}}\big) \un{J\,mol^{-1}}\\
\Eenr \equiv \frac{\cl}{\cb} &= \exp\!\Big(-\frac{\DGseq}{\Rg T}\Big) \un{1}
\end{align}
Valid for dilute Ca$^{2+}$ with unsaturated sites ($\Eenr\cb\ll\cmax$) and equal activity coefficients in bulk and layer (so $\Eenr$ is both a concentration and an activity ratio). Saturation (Langmuir closure):
\begin{gather}
\cl=\frac{\Eenr\cb}{1+\Eenr\cb/\cmax}\le\min(\Eenr\cb,\ \cmax) \un{mol\,m^{-3}}\\
\cmax=\Gs/\ds \un{mol\,m^{-3}}\nonumber
\end{gather}

\hd{3\quad Supersaturation (one definition)}
\phys{Crystallisation is driven by the excess chemical potential of the dissolved ion pair over the solid of the registered framing (ACC or calcite).}
\begin{equation}
\Om=\frac{\mathrm{IAP}}{\Ksp}=\frac{\aCa\,\aCO}{\Ksp},\qquad \Dmu=\kB T\ln\Om \un{J}
\end{equation}
per formula unit; nucleation requires $\Om>1$ ($\Dmu>0$). Only Ca$^{2+}$ is enriched; carbonate stays at its bulk activity:
\begin{gather}
\Oml=\frac{\cl}{\cb}\,\Omb \;\xrightarrow{\text{dilute}}\; \Eenr\,\Omb \un{1}\\
\Dmu_\mathrm{local}=\Dmu_\mathrm{bulk}-\DGseq/\NA \un{J}
\end{gather}

\hd{4\quad Classical nucleation (C3, \Rs)}
\phys{A cluster pays surface free energy before its volume gains $\Dmu$; the barrier is the top of that trade-off. On a substrate the cluster is a spherical cap with the same curvature.}
\begin{gather}
\DGhom=\frac{16\pi\,\gnl^{3}\vm^{2}}{3\,\Dmu^{2}},\quad r^{*}=\frac{2\gnl\vm}{\Dmu},\quad \vm=\frac{\Vm}{\NA} \un{J;\ m}\\
\DGhet=f(\theta)\,\DGhom,\qquad f(\theta)=\tfrac14(2+\cos\theta)(1-\cos\theta)^{2}
\end{gather}
$f(0)=0$, $f(\pi/2)=\tfrac12$, $f(\pi)=1$; $f'(\theta)=\tfrac34\sin^{3}\theta\ge0$, so $0\le\DGhet\le\DGhom$ on $[0,\pi]$. $\DGhom\to\infty$ as $\Om\to1^{+}$, and $\DGhom\to0$ as $\Om\to\infty$.

\hd{5\quad Young, with one sign convention}
\phys{$\theta$ is measured through the nucleus. Forming the contact removes substrate--liquid interface and creates substrate--nucleus interface over the same disc; $\cos\theta$ is that energy gain per unit nucleus--liquid energy.}
\begin{equation}
\cos\theta=\frac{\gsl-\gsn}{\gnl},\qquad \gsn=\gsnz+\dgw \un{1;\ J\,m^{-2}}
\end{equation}
$\theta=0$: the nucleus wets the template completely (no barrier); $\theta=\pi$: no wetting (homogeneous barrier).
$\dgw\equiv\gsn-\gsnz$ is the free energy per unit contact area to displace the template-bound water layer; \textbf{$\dgw>0$ is a cost} and raises $\theta$, $f$ and $\DGhet$ ($\partial\cos\theta/\partial\dgw=-1/\gnl<0$). If $\gsn$ is supplied as a full in-water interfacial free energy, $\dgw\equiv0$ and is not added again.
$\geff\equiv\gnl-\dgw$ (deck form) is \textbf{not} used: Young and $\DGhom$ take the bare $\gnl$. Reason: $\gnl$ multiplies the cap's free surface, which never touches the template, and $\DGhet=f\,\DGhom$ follows only if the same $\gnl$ sits in Young and in $\DGhom$. With $\geff$ in the denominator, $\partial\cos\theta/\partial\dgw=(\gsl-\gsnz)/\geff^{2}>0$ on the wetting side: a cost would lower the barrier.

\hd{6\quad Where the template acts}
\phys{The template enters the barrier through two channels only: wetting (interfacial energies, via $\theta$) and Ca$^{2+}$ enrichment (via $\Om$).}
\begin{equation}
\DGhet=f\big(\theta[\gsl,\gsnz,\dgw]\big)\;\frac{16\pi\gnl^{3}\vm^{2}}{3\big(\kB T\ln\Oml[\DGseq,\Gs,\ds]\big)^{2}}
\end{equation}
\textbf{No-double-counting condition:} each template-dependent input enters exactly one channel ($\partial\theta/\partial\DGseq=0$, $\partial\Om/\partial\gamma_\mathrm{s\cdot}=0$), and $\gsl,\gsn$ are evaluated without the template-bound Ca$^{2+}$ excess already credited to $\Oml$. Bulk inputs ($\gnl,\Vm,T$) may enter both. Proof: P0-W3.2.

\hd{7\quad Conformality (C4, \Rs)}
\phys{A barrier needs a continuous film; whether the mineral spreads or beads is a wetting balance, not a barrier height.}
\begin{equation}
\Ssp=\gsl-\gsn-\gnl=\gnl(\cos\theta-1) \un{J\,m^{-2}}
\end{equation}
$\Ssp\ge0$: no finite $\theta$ ($\theta\to0$, $f\to0$), layer-by-layer growth. $\Ssp<0$: $\theta>0$, island growth; a film forms only on coalescence, at cap height $\hc=\tan(\theta/2)/(2\sqrt{\Ns})$ [m]. A low $f$ does not imply coverage: $f(\pi/2)=\tfrac12$ is still island growth.

\hd{8\quad Stability (C5): screens, necessary not sufficient}
(a) \Rs, aqueous Pourbaix, per non-O, non-H atom [J\,mol$^{-1}$]:
\begin{multline}
\Psi_\phi=\tfrac{1}{n_\mathrm{M}}\big[G_\phi-n_\mathrm{O}\mu_{\mathrm{H_2O}}-\Rg T\ln10\,(2n_\mathrm{O}-n_\mathrm{H})\,\mathrm{pH}\\
-(2n_\mathrm{O}-n_\mathrm{H}+Q)\Fc\ESHE\big]
\end{multline}
$\DGpbx=\Psi_\phi-\Psi_\mathrm{hull}\ge0$; zero means stable at (pH, $\ESHE$).

(b) \Rc, Li open, $\ELi$ vs Li$^+$/Li [J\,mol$^{-1}$; V]:
\begin{equation}
\muLi=\muLiz-\Fc\ELi,\qquad \Phi_\phi=G_\phi-n_\mathrm{Li}\muLi
\end{equation}
Window $[E_\mathrm{red},E_\mathrm{ox}]$: the $\ELi$ over which $\phi$ lies on the Li-open hull.

\hd{9\quad Insulating film (C4 limitation, \Rc)}
\phys{The mineral is an electronic insulator, so conformality cannot be bought with thickness; the film must still pass the collector current.}
\begin{gather}
R_A=\rhoe d\le R_{A,\max}\;\Rightarrow\; d\le R_{A,\max}/\rhoe \un{\Omega\,m^{2};\ m}\\
j=j_0\e^{-2\kappa d},\ \ \kappa=\frac{\sqrt{2m_\mathrm{e}\phiB}}{\hbar}\;\Rightarrow\; d\le\frac{\ln(j_0/\jreq)}{2\kappa} \un{m}
\end{gather}
\end{multicols}

\newpage
{\large\bfseries Symbols} \hfill {\small generated from \texttt{theory/symbols.yaml}; Reg.: S = \Rs, C = \Rc; -- = dimensionless}

\smallskip
% Generated by theory/tools/gen_symbol_table.py from theory/symbols.yaml; do not edit.
% Regime column: S = R-slurry, C = R-cell.
\begin{minipage}[t]{\dimexpr(\textwidth-2mm)/2\relax}\footnotesize\setlength{\tabcolsep}{2pt}\begin{tabular}[t]{@{}l p{\dimexpr\linewidth*29/50\relax} l l@{}}\hline
Symbol & Meaning & SI unit & Reg.\\\hline
\multicolumn{4}{@{}l}{\textit{Exact SI constants (kept symbolic)}}\\
$k_\mathrm{B}$ & Boltzmann constant & $\mathrm{J}\,\mathrm{K}^{-1}$ & S,C\\
$N_\mathrm{A}$ & Avogadro constant & $\mathrm{mol}^{-1}$ & S,C\\
$R$ & molar gas constant & $\mathrm{J}\,\mathrm{K}^{-1}\,\mathrm{mol}^{-1}$ & S,C\\
$F$ & Faraday constant & $\mathrm{C}\,\mathrm{mol}^{-1}$ & S,C\\
$e$ & elementary charge & $\mathrm{C}$ & S,C\\
$\hbar$ & reduced Planck constant & $\mathrm{J}\,\mathrm{s}$ & C\\
$m_\mathrm{e}$ & electron mass & $\mathrm{kg}$ & C\\
\multicolumn{4}{@{}l}{\textit{State variables}}\\
$T$ & absolute temperature of the regime & $\mathrm{K}$ & S,C\\
$\mathrm{pH}$ & pH (molal scale) & -- & S\\
$E_\mathrm{SHE}$ & collector potential vs SHE & $\mathrm{V}$ & S\\
$E_\mathrm{Li}$ & collector potential vs Li\textsuperscript{+}/Li & $\mathrm{V}$ & C\\
\multicolumn{4}{@{}l}{\textit{Adhesion in water (C1)}}\\
$\gamma_\mathrm{ow}$ & oxide-water interfacial free energy & $\mathrm{J}\,\mathrm{m}^{-2}$ & S\\
$\gamma_\mathrm{pw}$ & polymer-water interfacial free energy & $\mathrm{J}\,\mathrm{m}^{-2}$ & S\\
$\gamma_\mathrm{op}$ & oxide-polymer interfacial free energy in contact & $\mathrm{J}\,\mathrm{m}^{-2}$ & S\\
$\gamma_\mathrm{o}$ & oxide surface free energy against vacuum & $\mathrm{J}\,\mathrm{m}^{-2}$ & S\\
$\gamma_\mathrm{p}$ & polymer surface free energy against vacuum & $\mathrm{J}\,\mathrm{m}^{-2}$ & S\\
$\gamma_\mathrm{w}$ & water surface free energy against its vapour & $\mathrm{J}\,\mathrm{m}^{-2}$ & S\\
$\gamma_\mathrm{n}$ & nucleus surface free energy against vacuum (dry) & $\mathrm{J}\,\mathrm{m}^{-2}$ & S\\
$W_\mathrm{adh,wet}$ & work of adhesion in water (Dupre) & $\mathrm{J}\,\mathrm{m}^{-2}$ & S\\
$W_\mathrm{adh,dry}$ & work of adhesion against vacuum (comparator) & $\mathrm{J}\,\mathrm{m}^{-2}$ & S\\
$\Delta F_\mathrm{PMF}$ & PMF free energy, separated minus bound & $\mathrm{J}$ & S\\
$A_\mathrm{c}$ & fragment contact area & $\mathrm{m}^{2}$ & S\\
\multicolumn{4}{@{}l}{\textit{Sequestration (C2)}}\\
$\Delta G_\mathrm{seq}$ & Ca\textsuperscript{2+} transfer free energy, bulk to binding layer & $\mathrm{J}\,\mathrm{mol}^{-1}$ & S\\
$\Delta G_\mathrm{site}$ & transfer free energy to the deprotonated site & $\mathrm{J}\,\mathrm{mol}^{-1}$ & S\\
$\mathrm{p}K_\mathrm{a}$ & pKa of the coordinating group & -- & S\\
$\mathcal{E}$ & Ca\textsuperscript{2+} enrichment factor (dilute limit) & -- & S\\
$c_\mathrm{bulk}$ & bulk Ca\textsuperscript{2+} concentration & $\mathrm{mol}\,\mathrm{m}^{-3}$ & S\\
$c_\mathrm{local}$ & mean Ca\textsuperscript{2+} concentration in the layer & $\mathrm{mol}\,\mathrm{m}^{-3}$ & S\\
$c_\mathrm{max}$ & saturation concentration of the layer & $\mathrm{mol}\,\mathrm{m}^{-3}$ & S\\
$\Gamma_\mathrm{s}$ & areal binding-site density & $\mathrm{mol}\,\mathrm{m}^{-2}$ & S\\
$\delta_\mathrm{s}$ & binding-layer thickness & $\mathrm{m}$ & S\\
$V_\mathrm{site}$ & volume per binding site & $\mathrm{m}^{3}$ & S\\
$a_\mathrm{site}$ & footprint area of one binding site & $\mathrm{m}^{2}$ & S\\
$\rho_\mathrm{L}$ & volume density of coordinating groups & $\mathrm{mol}\,\mathrm{m}^{-3}$ & S\\
$d_\mathrm{layer}$ & template layer thickness & $\mathrm{m}$ & S\\
$r_\mathrm{Ca}$ & Ca\textsuperscript{2+} ionic radius & $\mathrm{m}$ & S\\
$x$ & Langmuir binding variable & -- & S\\
$\Theta$ & fraction of sites occupied by Ca\textsuperscript{2+} & -- & S\\
$\Gamma_\mathrm{Ca}$ & Ca\textsuperscript{2+} surface excess (Gibbs) & $\mathrm{mol}\,\mathrm{m}^{-2}$ & S\\
$\Gamma_\mathrm{Ca,sn}$ & Ca\textsuperscript{2+} excess at the template-nucleus contact & $\mathrm{mol}\,\mathrm{m}^{-2}$ & S\\
\multicolumn{4}{@{}l}{\textit{Carbonate speciation (activities)}}\\
$m^{\circ}$ & standard molality & $\mathrm{mol}\,\mathrm{kg}^{-1}$ & S\\
$\rho_\mathrm{w}$ & density of water & $\mathrm{kg}\,\mathrm{m}^{-3}$ & S\\
$I$ & ionic strength (molal) & $\mathrm{mol}\,\mathrm{kg}^{-1}$ & S\\
$C_\mathrm{T}$ & dissolved inorganic carbon & $\mathrm{mol}\,\mathrm{m}^{-3}$ & S\\
$\mathrm{Ca}_\mathrm{T}$ & total dissolved calcium & $\mathrm{mol}\,\mathrm{m}^{-3}$ & S\\
$\mathrm{p}K_1$ & pK1 of carbonic acid & -- & S\\
$\mathrm{p}K_2$ & pK2 of carbonic acid & -- & S\\
$K_{\mathrm{CaCO_3^0}}$ & formation constant of CaCO\textsubscript{3}(aq) & -- & S\\
$K_{\mathrm{CaHCO_3^+}}$ & formation constant of CaHCO3+ & -- & S\\
${}^{*}K_{\mathrm{CaOH^+}}$ & hydrolysis constant of CaOH+ & -- & S\\
$A_\mathrm{DH}$ & Debye-Hueckel constant & -- & S\\
$b_\mathrm{D}$ & Davies linear coefficient & -- & S\\
$I_\mathrm{val}$ & upper validity bound of the activity model & $\mathrm{mol}\,\mathrm{kg}^{-1}$ & S\\
$\gamma_{\pm1}$ & activity coefficient, charge one & -- & S\\
$\gamma_{\pm2}$ & activity coefficient, charge two & -- & S\\
$m_{C_\mathrm{T}}$ & dissolved inorganic carbon, molal & -- & S\\
$m_{\mathrm{Ca}_\mathrm{T}}$ & total calcium, molal & -- & S\\
$A_\mathrm{c}$ & carbonate speciation factor & -- & S\\
$P_\mathrm{c}$ & ion-pair factor & -- & S\\
$D_\mathrm{c}$ & free-calcium factor & -- & S\\
\multicolumn{4}{@{}l}{\textit{Supersaturation}}\\
$a_{\mathrm{Ca}^{2+}}$ & bulk Ca\textsuperscript{2+} activity & -- & S\\
$a_{\mathrm{CO}_3^{2-}}$ & bulk carbonate activity & -- & S\\
$K_\mathrm{sp}$ & solubility product of the registered framing & -- & S\\
$a_\mathrm{w}$ & water activity of the regime state & -- & S\\
$n$ & structural water per formula unit (hydrated framing) & -- & S\\
$\Omega_\mathrm{bulk}$ & bulk saturation ratio IAP/Ksp & -- & S\\
$\Omega$ & saturation ratio where the nucleus forms & -- & S\\
$\Delta\mu$ & driving force per formula unit & $\mathrm{J}$ & S\\
\hline\end{tabular}\end{minipage}\hfill
\begin{minipage}[t]{\dimexpr(\textwidth-2mm)/2\relax}\footnotesize\setlength{\tabcolsep}{2pt}\begin{tabular}[t]{@{}l p{\dimexpr\linewidth*29/50\relax} l l@{}}\hline
Symbol & Meaning & SI unit & Reg.\\\hline
\multicolumn{4}{@{}l}{\textit{Nucleation and Young (C3)}}\\
$\gamma_\mathrm{nl}$ & nucleus-liquid interfacial free energy & $\mathrm{J}\,\mathrm{m}^{-2}$ & S\\
$\gamma_\mathrm{sl}$ & template-liquid energy at the regime activities & $\mathrm{J}\,\mathrm{m}^{-2}$ & S\\
$\gamma_\mathrm{sl}^{(0)}$ & template-liquid energy without Ca\textsuperscript{2+} binding & $\mathrm{J}\,\mathrm{m}^{-2}$ & S\\
$\Delta\gamma_\mathrm{ads}$ & Ca\textsuperscript{2+} adsorption change of the template-liquid energy & $\mathrm{J}\,\mathrm{m}^{-2}$ & S\\
$\gamma_\mathrm{sn}$ & template-nucleus energy in water (= ref + correction) & $\mathrm{J}\,\mathrm{m}^{-2}$ & S\\
$\gamma_\mathrm{sn}^{(0)}$ & template-nucleus energy, reference without water displacement & $\mathrm{J}\,\mathrm{m}^{-2}$ & S\\
$\Delta\gamma_\mathrm{w}$ & water-displacement correction (cost $>$ 0) & $\mathrm{J}\,\mathrm{m}^{-2}$ & S\\
$\gamma_\mathrm{eff}$ & deck effective nucleus-liquid energy (audit only) & $\mathrm{J}\,\mathrm{m}^{-2}$ & S\\
$\theta$ & contact angle, measured through the nucleus & $\mathrm{rad}$ & S\\
$f(\theta)$ & heterogeneous shape factor & -- & S\\
$V_\mathrm{m}$ & molar volume of the registered framing & $\mathrm{m}^{3}\,\mathrm{mol}^{-1}$ & S\\
$r^{*}$ & critical radius & $\mathrm{m}$ & S\\
$\Delta G^{*}_\mathrm{hom}$ & homogeneous nucleation barrier per nucleus & $\mathrm{J}$ & S\\
$\Delta G^{*}_\mathrm{het}$ & heterogeneous barrier on the template & $\mathrm{J}$ & S\\
\multicolumn{4}{@{}l}{\textit{Accessibility of the barrier (C3)}}\\
$J_{0}$ & heterogeneous nucleation prefactor per area & $\mathrm{m}^{-2}\,\mathrm{s}^{-1}$ & S\\
$t_\mathrm{p}$ & time available for nucleation & $\mathrm{s}$ & S\\
$N_\mathrm{req}$ & required areal density of nuclei & $\mathrm{m}^{-2}$ & S\\
$q_0$ & admissible fraction of empty contact patches & -- & S\\
$A_\mathrm{patch}$ & particle-collector contact patch & $\mathrm{m}^{2}$ & S\\
$B_\mathrm{acc}$ & largest accessible barrier & $\mathrm{J}$ & S\\
$N_\mathrm{nuc}$ & expected areal density of nuclei & $\mathrm{m}^{-2}$ & S\\
\multicolumn{4}{@{}l}{\textit{Conformality (C4)}}\\
$S_\mathrm{spread}$ & spreading coefficient & $\mathrm{J}\,\mathrm{m}^{-2}$ & S\\
$N_\mathrm{s}$ & island density at coalescence & $\mathrm{m}^{-2}$ & S\\
$h_\mathrm{c}$ & geometric coalescence thickness & $\mathrm{m}$ & S\\
$h_\mathrm{cl}$ & geometric closure thickness & $\mathrm{m}$ & S\\
$h_\mathrm{coal}$ & kinetic critical coalescence thickness (charter C4 (c)) & $\mathrm{m}$ & S\\
\multicolumn{4}{@{}l}{\textit{Stability (C5)}}\\
$G_\phi$ & molar Gibbs energy of species phi & $\mathrm{J}\,\mathrm{mol}^{-1}$ & S\\
$\mu_{\mathrm{H_2O}}$ & chemical potential of water & $\mathrm{J}\,\mathrm{mol}^{-1}$ & S\\
$n_\mathrm{O}$ & number of O atoms per formula unit & -- & S\\
$n_\mathrm{H}$ & number of H atoms per formula unit & -- & S\\
$n_\mathrm{M}$ & number of non-O, non-H atoms & -- & S\\
$Q$ & species charge in units of e & -- & S\\
$\Psi_\phi$ & Pourbaix potential per metal atom & $\mathrm{J}\,\mathrm{mol}^{-1}$ & S\\
$\Psi_\mathrm{hull}$ & lower envelope of Psi (the hull) & $\mathrm{J}\,\mathrm{mol}^{-1}$ & S\\
$\Delta G_\mathrm{pbx}$ & energy above the aqueous Pourbaix hull & $\mathrm{J}\,\mathrm{mol}^{-1}$ & S\\
$n_\mathrm{Li}$ & number of Li atoms per formula unit & -- & C\\
$\mu_\mathrm{Li}^{\circ}$ & Li chemical potential in Li metal & $\mathrm{J}\,\mathrm{mol}^{-1}$ & C\\
$\mu_\mathrm{Li}$ & Li chemical potential at the electrode potential & $\mathrm{J}\,\mathrm{mol}^{-1}$ & C\\
$\Phi_\phi$ & Li-open grand potential & $\mathrm{J}\,\mathrm{mol}^{-1}$ & C\\
$\Delta G_\mathrm{r}$ & reaction free energy against Li metal & $\mathrm{J}\,\mathrm{mol}^{-1}$ & C\\
$\Delta n$ & Li released by the reaction & -- & C\\
$\Delta\Phi$ & grand-potential change of the reaction & $\mathrm{J}\,\mathrm{mol}^{-1}$ & C\\
$E^{*}$ & potential limit set by one reaction & $\mathrm{V}$ & C\\
$E_\mathrm{red}$ & lower limit of the stability window & $\mathrm{V}$ & C\\
$E_\mathrm{ox}$ & upper limit of the stability window & $\mathrm{V}$ & C\\
\multicolumn{4}{@{}l}{\textit{Insulating film (C4, R-cell)}}\\
$d$ & film thickness & $\mathrm{m}$ & C\\
$\rho_\mathrm{e}$ & film electronic resistivity & $\Omega\,\mathrm{m}$ & C\\
$R_A$ & area-specific film resistance & $\Omega\,\mathrm{m}^{2}$ & C\\
$R_{A,\max}$ & admissible area-specific resistance (input) & $\Omega\,\mathrm{m}^{2}$ & C\\
$\phi_\mathrm{B}$ & tunnelling barrier height & $\mathrm{J}$ & C\\
$\kappa$ & tunnelling decay constant & $\mathrm{m}^{-1}$ & C\\
$j_0$ & tunnelling prefactor current density & $\mathrm{A}\,\mathrm{m}^{-2}$ & C\\
$j_\mathrm{req}$ & required current density (input) & $\mathrm{A}\,\mathrm{m}^{-2}$ & C\\
$d_\mathrm{max}^\mathrm{tun}$ & thickness bound, direct tunnelling & $\mathrm{m}$ & C\\
$d_\mathrm{max}$ & admissible film thickness & $\mathrm{m}$ & C\\
$d_\mathrm{max}^\mathrm{ohm}$ & thickness bound, ohmic limit & $\mathrm{m}$ & C\\
\multicolumn{4}{@{}l}{\textit{Registration constants (theory/thresholds.md)}}\\
$p$ & credible level of the decision rule & -- & S,C\\
$p_\mathrm{deg}$ & least mass shift that counts as degradation & -- & S,C\\
$m_W$ & C1 margin on the wet work of adhesion & $\mathrm{J}\,\mathrm{m}^{-2}$ & S\\
$m_\mathrm{seq}$ & C2 margin on the Ca\textsuperscript{2+} transfer free energy & $\mathrm{J}\,\mathrm{mol}^{-1}$ & S\\
$m_{\cos}$ & C3 margin on cos theta & -- & S\\
$m_S$ & C4 margin on the spreading coefficient & $\mathrm{J}\,\mathrm{m}^{-2}$ & S\\
$e_\mathrm{f}$ & expected error of a computed formation energy & $\mathrm{J}\,\mathrm{mol}^{-1}$ & S,C\\
$n_\mathrm{at}$ & atoms per formula unit & -- & S,C\\
$\Delta_\mathrm{ACC}$ & ACC offset above calcite & $\mathrm{J}\,\mathrm{mol}^{-1}$ & S\\
$\tau_\mathrm{hull}$ & C5 hull tolerance & $\mathrm{J}\,\mathrm{mol}^{-1}$ & S\\
$m_E$ & C5 margin on the window limits & $\mathrm{V}$ & C\\
$m_\mathrm{expt}$ & C5 experimental margin on the window limits & $\mathrm{V}$ & C\\
$E_\mathrm{cell}^{\min}$ & lower end of the R-cell potential range & $\mathrm{V}$ & C\\
$E_\mathrm{cell}^{\max}$ & upper end of the R-cell potential range & $\mathrm{V}$ & C\\
\hline\end{tabular}\end{minipage}


\smallskip
{\small \textbf{Conventions.} Angles in rad; free energies per formula unit (J) or per mole (J\,mol$^{-1}$) as stated; activities on the molal scale; potentials vs SHE in \Rs\ and vs Li$^+$/Li in \Rc. $\Om$ is the only supersaturation variable (the charter's $S$ for C3 is $\Om$). The functions W4 implements, with exact limits, monotonicity and the sign-flip tests, are in \texttt{theory/contract.yaml}.}
\end{document}
